\documentclass{exam}

\newif\ifA
\newif\ifkey
\newif\ifprintselected   % required by exam_config's \ansbox

%%%%%%%%%%%%%%%%%%%%%%%%%%%%%%%%%%%%%%%%%%%%%%%%%%%%%%
%%% SET THESE IF VARIABLES
%%%%%%%%%%%%%%%%%%%%%%%%%%%%%%%%%%%%%%%%%%%%%%%%%%%%%%
% Lab Num -- Module 3's six lab periods follow L04 (Lab 1) and the two
% Module 2 measurement labs, so L28 is Lab 9. Override with \def\labNumIn{}
% if the course sequence renumbers.
\ifdefined\labNumIn
	\def\labNum{\labNumIn}
\else
	\def\labNum{9}
\fi

% is this the key or not?
\ifdefined\iskey
	\ifnum\iskey=1
		\keytrue
	\else
		\keyfalse
	\fi
\else
	%% Pick one manually
	\keyfalse % if this is NOT the key
%	\keytrue     % if this IS the key
\fi

% set up the header and footer - change for every term
\ifdefined\thisTermIn
	\def\thisterm{\thisTermIn}
\else
	\def\thisterm{Fall 2026}
\fi
%
\def\thisclass{ECE 444}

%%%%%%%%%%%%%%%%%%%%%%%%%%%%%%%%%%%%%%%%%%%%%%%%%%%%%%
%%% END SET THESE IF VARIABLES
%%%%%%%%%%%%%%%%%%%%%%%%%%%%%%%%%%%%%%%%%%%%%%%%%%%%%%
% packages and shortcuts
\input{myPackages_gr}
\usepackage[hidelinks, linkcolor=red]{hyperref}
\input{myShortcuts}
\setlength{\parindent}{0pt}
\setlength{\parskip}{6pt}

%% exam class customization
\input{exam_config}
\printselectedfalse      % full worked solutions in the key, not answers-only
\CorrectChoiceEmphasis{\color{red}}
\SolutionEmphasis{\color{red}}
\renewcommand\questionlabel{\textbf{\thequestion.}}
\renewcommand{\thepartno}{\alph{partno}}
\renewcommand{\partlabel}{(\thepartno)}

% Fillable table cell: blank in the student copy, the expected value in red on
% the KEY. Fixed-width centered columns keep the blank cells writable.
\newcolumntype{K}[1]{>{\centering\arraybackslash}p{#1}}
\newcommand{\kv}[1]{\ifkey{\color{red}#1}\fi\rule[-0.10in]{0pt}{0.34in}}

% print the answers if this is the key
\ifkey
	\printanswers % if this IS the key
\else
	\noprintanswers  %if this is NOT the key
\fi

% print the header and footer
% need to print key on top if this is the key
\ifkey
	\firstpageheader{}{\color{red}{ KEY}}{}
	\runningheader{\thisclass\hspace{1pt} -- \thisterm }{Lab \#\labNum\hspace{1pt}\color{red}{ KEY}}{\thepage\hspace{1pt} of\hspace{1pt} \numpages}
	\firstpagefooter{}{}{}
	\runningfooter{}{}{}
\else
	\firstpageheader{}{}{}
	\runningheader{\thisclass \hspace{1pt} -- \thisterm }{Lab \#\labNum\hspace{1pt}}{\thepage\hspace{1pt} of\hspace{1pt} \numpages}
	\firstpagefooter{}{}{}
	\runningfooter{}{}{}
\fi

\runningheadrule

\begin{document}
\pagestyle{headandfoot}   % exam class


%%%%%%%%%%%%%%%%%%%%%%%%%%%%%%%%%%%%%%%%%%%%%%%%%%%%%%
% title block
%%%%%%%%%%%%%%%%%%%%%%%%%%%%%%%%%%%%%%%%%%%%%%%%%%%%%%

\begin{center}
USAF ACADEMY \\
DEPARTMENT OF ELECTRICAL AND COMPUTER ENGINEERING \\
LAB \#\labNum\ (Lesson 28) - Null Steering Lab \\
\textbf{Lab Sheet} \\
\thisterm \\[4mm]
\end{center}

\vspace{-4pt}
\textbf{Name:} \rule[-2pt]{2.6in}{0.4pt} \hfill
\textbf{Section:} \rule[-2pt]{0.9in}{0.4pt} \hfill
\textbf{Date:} \rule[-2pt]{1.1in}{0.4pt}

\vspace{4pt}

\fullwidth{\textbf{Learning Objective 3.9: } Calculate null-steering weights,
implement pattern nulls on the ADALM-PHASER, and compare manual null steering
with adaptive beamforming.}

\vspace{2pt}

This sheet is the turn-in document for the Lesson 28 lab. Every measurement
below is a comparison against a frozen reference trace, so record the reference
before you change anything. The lab is a team assignment, so one sheet per team.

\begin{questions}

%%%%%%%%%%%%%%%%%%%%%%%%%%%%%%%%%%%%%%%%%%%%%%%%%%%
\question \textbf{The static notch.} Compare the null-steered sweep against the
frozen uniform reference. The levels at $+22.5\mydeg$ are relative to their own
trace's peak, and the main-lobe entries are relative to the uniform peak.

\begin{center}
\setlength{\tabcolsep}{4pt}
\renewcommand{\arraystretch}{1.25}
\begin{tabular}{| l | K{1.6in} | K{1.7in} |}
\hline
\textbf{Quantity} & \textbf{Measured} & \textbf{Predicted} \\\hline\hline
Reference level at $+22.5\mydeg$ & \kv{$-11$ to $-13$ dBc} & \kv{$-12.8$ dBc} \\\hline
Null-steered level at $+22.5\mydeg$ & \kv{$-20$ to $-22$ dBc} & \kv{$-21.6$ dBc} \\\hline
Change at that angle & \kv{$8$ to $10$ dB deeper} & \kv{$8.8$ dB deeper} \\\hline
Main-lobe peak, uniform & \kv{reference} & \kv{$0$ dB} \\\hline
Main-lobe peak, null-steered & \kv{$1.5$ to $2.5$ dB down} & \kv{$1.8$ dB down} \\\hline
Main-lobe cost & \kv{$1.5$ to $2.5$ dB} & \kv{$1.8$ dB} \\\hline
\end{tabular}
\end{center}

\begin{solution}
The weights come from $w = w_d - r_n w_n$ with
$r_n = (w_n^H w_d)/(w_n^H w_n)$, entered as the eight gains
75, 65, 82, 100, 100, 82, 65, and 75 percent and the eight phase offsets
$-12.1\mydeg$, $+3.1\mydeg$, $+13.0\mydeg$, $+6.0\mydeg$, $-6.0\mydeg$,
$-13.0\mydeg$, $-3.1\mydeg$, $+12.1\mydeg$. The exact main-lobe cost is $2.0$ dB
and the sweep reads about $1.8$ dB. The notch bottoms out near $-21.6$ dBc
rather than going to zero: the floor sits about 23 dB below the uniform peak and
the null weights cost about 2 dB of main lobe, so roughly 21 dB below the
null-steered peak is the deepest notch this plot can report. A sign error on the
phase column moves the notch to $-22.5\mydeg$, which is the most common failure
here.
\end{solution}

%%%%%%%%%%%%%%%%%%%%%%%%%%%%%%%%%%%%%%%%%%%%%%%%%%%
\newpage
\question \textbf{The difference null.} With Beam 1 Phase set to $180\mydeg$,
record the depth of the boresight null and the angles of the two peaks that
replace the main lobe.

\begin{center}
\setlength{\tabcolsep}{4pt}
\renewcommand{\arraystretch}{1.25}
\begin{tabular}{| l | K{1.7in} | K{1.6in} |}
\hline
\textbf{Quantity} & \textbf{Measured} & \textbf{Predicted} \\\hline\hline
Null depth at boresight & \kv{$-20$ to $-23$ dBc} & \kv{$-22$ dBc} \\\hline
Left peak angle & \kv{$-11\mydeg \pm 2.8\mydeg$} & \kv{$-11\mydeg$} \\\hline
Right peak angle & \kv{$+11\mydeg \pm 2.8\mydeg$} & \kv{$+11\mydeg$} \\\hline
\end{tabular}
\end{center}

\begin{solution}
Nothing in the analog beamformers changed for this measurement. All eight
elements still carry the phases they had a moment earlier, and the only change
is a sign on one of the two digital channels, applied after the signal was
digitized. That is enough to cancel the two subarray outputs against each other
on boresight. The peak angles are read on the same $2.8125\mydeg$ grid as every
other sweep, so one grid step is the tolerance. This trace is the monopulse
difference beam, and the Tracking run in Part 7 shows the same null at
$-21.8$ dBc.
\end{solution}

%%%%%%%%%%%%%%%%%%%%%%%%%%%%%%%%%%%%%%%%%%%%%%%%%%%
\question \textbf{MVDR against manual.} From your own procedure C sweeps, with
the second HB100 at roughly $+30\mydeg$ and running about 10 dB stronger than
the boresight source. Report both levels relative to the peak of the manual
trace.

\begin{center}
\setlength{\tabcolsep}{4pt}
\renewcommand{\arraystretch}{1.25}
\begin{tabular}{| l | K{1.9in} | K{1.9in} |}
\hline
\textbf{Mode} & \textbf{Level at the interferer angle} & \textbf{Level in the look direction} \\\hline\hline
Manual & \kv{$0$ dB, the trace peaks there} & \kv{$\approx -10$ dB} \\\hline
MVDR & \kv{$-17$ to $-19$ dB} & \kv{within $\approx 1$ dB of the manual reading} \\\hline
Difference & \kv{$17$ to $19$ dB} & \kv{$\le 1$ dB} \\\hline
\end{tabular}
\end{center}

\begin{solution}
The manual beamformer is captured by the interferer. With 10 dB in its favor the
second source dominates the received power, so the manual trace peaks toward it
and the boresight target sits about 10 dB down. MVDR estimates
$\hat R = \frac{1}{K}XX^H$ over $K = 128$ snapshots and solves
$w = R^{-1}s/(s^H R^{-1} s)$, which holds unit gain in the look direction while
minimizing total output power, so the look-direction level survives and the
response toward the interferer drops by 17 to 19 dB. The measurement is the gap
between the two traces read at the interferer's angle. The FFT tab cannot
separate the two sources, because both HB100s sit at nominally the same
frequency, so the evidence is the shape of the sweep.
\end{solution}

%%%%%%%%%%%%%%%%%%%%%%%%%%%%%%%%%%%%%%%%%%%%%%%%%%%
\newpage
\question \textbf{Written answers.} Three or four sentences each.

\begin{parts}

	\part Why is the measured notch depth limited by the sweep's noise floor
	rather than by quantization? Give the numbers that set the floor, and the
	depth the quantized weights alone would allow.

	\begin{solutionorlines}[2.4in]
	The floor of the sweep sits about 23 dB below the uniform-taper peak, and the
	null weights cost about 2 dB of main lobe, so the deepest notch the plot can
	report is roughly 21 dB below the null-steered peak. That is the $20$ to
	$22$ dB recorded above. Quantization is not the limit: rounding each phase to
	the $2.8125\mydeg$ LSB and each gain to its 1 percent step produces a slightly
	different weight vector, which nulls a direction a fraction of a degree away
	from the one requested, and pushing those quantized weights through the array
	factor leaves about $-48$ dB at the designed angle. That residual is more than
	25 dB below anything this measurement can display, so the plot is reporting
	the dynamic range of the measurement rather than the resolution of the phase
	shifter.
	\end{solutionorlines}

	\part Give one situation where a computed static null is the better choice
	than MVDR, and one where it is not, with a reason for each.

	\begin{solutionorlines}[2.4in]
	A computed static null is the better choice when the interference sits in a
	known, fixed direction, such as a stationary emitter or a known clutter
	bearing. It uses all eight analog weights, so it can place several nulls at
	once and shape the rest of the pattern deliberately, and it needs no snapshots
	and no digital channels. It is the wrong choice when the interferer is unknown
	or moving, because every change of geometry means recomputing eight weights
	and entering them again, and the null is only as good as the angle that was
	assumed. MVDR finds its own null from the received data several times a
	second, at the cost of the two digital degrees of freedom this board has,
	which is one constraint and one null.
	\end{solutionorlines}

\end{parts}

\end{questions}

\vspace{1.0em}
\noindent\textbf{Documentation:} \rule[-2pt]{0.6\linewidth}{0.4pt}

\end{document}
